\documentclass[a4paper,12pt]{jsarticle}

% 明朝体指定（標準）
\usepackage{otf}

% 図表やコードを扱うためのパッケージ
\usepackage{graphicx}
\usepackage{verbatim}
\usepackage{amsmath}
\usepackage{listings}
\usepackage{array}

% ページ番号を下中央
\pagestyle{plain}

\begin{document}

%%%%%%%%%%%%%%%%%%%%%%%%%%%%%%%%%%%%%%%%%%%%
% 表紙
%%%%%%%%%%%%%%%%%%%%%%%%%%%%%%%%%%%%%%%%%%%%
\begin{titlepage}
\centering

\vspace*{4cm}

{\Large プログラミングB レポート \par}
\vspace{1cm}

{\Large 課題２（再帰下降構文解析による電卓）\par}

\vspace{3cm}

提出年月日：令和 7 年 12 月 xx 日 \\
教員名：武政 淳二 先生 \\[1cm]

提出者：\\
学生番号：XXXXXXX \\
氏名：山田 太郎

\end{titlepage}
\setcounter{page}{1}

%%%%%%%%%%%%%%%%%%%%%%%%%%%%%%%%%%%%%%%%%%%%
\section{課題内容}
%%%%%%%%%%%%%%%%%%%%%%%%%%%%%%%%%%%%%%%%%%%%

本レポートでは，「課題２」に取り組んだ。
具体的には，中置記法で与えられた算術式に対して，
再帰下降構文解析を用いて四則演算・括弧・指数演算を解析し，
値を計算する電卓プログラムを作成した。%
発展課題として，小数の扱い，指数演算，一元負号の処理，
および文法エラー検出を実装した。

%%%%%%%%%%%%%%%%%%%%%%%%%%%%%%%%%%%%%%%%%%%%
\section{アルゴリズムの説明}
%%%%%%%%%%%%%%%%%%%%%%%%%%%%%%%%%%%%%%%%%%%%

本課題では，中置記法の式を次の文法に基づいて解析する
再帰下降構文解析法を用いた。

\begin{itemize}
    \item $expression ::= term \ \{ (+ \mid -)\ term \}$
    \item $term ::= factor\ \{ (* \mid /)\ factor \}$
    \item $factor ::= constant \mid (expression) \mid -factor \mid factor\ ^\ digit$
    \item $constant ::= digits\ [.\ digits]$
\end{itemize}

式を「項」，項を「因子」に分解し，
各規則を関数に対応させ，
左から順に入力文字を読み進めながら再帰的に評価する。

アルゴリズムの要点は次の通りである。

\begin{itemize}
    \item 全体の計算は $expression$ を入口とする。
    \item $term$ が乗除を，$factor$ が括弧・指数・一元負号を処理する。
    \item 文字列の走査には 1 文字の前読み（lookahead）を用い，
    関数 \verb|next()| により現在文字 \verb|current_char| を更新する。
    \item 文法に合致しない文字を検出した場合は，入力位置とともに
    エラーを出力して終了する。
\end{itemize}

%%%%%%%%%%%%%%%%%%%%%%%%%%%%%%%%%%%%%%%%%%%%
\section{プログラムの説明}
%%%%%%%%%%%%%%%%%%%%%%%%%%%%%%%%%%%%%%%%%%%%

\subsection{入力形式}
本プログラムは，標準入力から 1 行の算術式を読み込む。
空白の混入は想定しない。

\subsection{出力形式}
式が正しい場合は計算結果を実数で出力し，
除算におけるゼロ除算では
\verb|Zero division occurs.| と表示し終了する。
文法エラーの場合は，
\verb|Syntax error occurs at the X-th character.| と出力する。

\subsection{データ構造}
入力式は文字配列 \verb|input[]| に格納し，
現在位置を整数 \verb|char_index|，
現在の文字を \verb|current_char| として保持する。

\subsection{大域変数}
\begin{verbatim}
char input[MAX_EXPRESSION_LENGTH];
int char_index;
char current_char;
\end{verbatim}

\subsection{関数の設計}

関数間の呼び出し関係は図\ref{fig:callgraph}の通りである。

\begin{figure}[h]
\centering
\begin{verbatim}
                    +-------------+
                    | expression  |
                    +------+------+ 
                           |
                           v
                    +-------------+
                    |    term     |
                    +------+------+ 
                           |
                           v
                    +-------------+
                    |   factor    |
                    +------+------+ 
                           |
                           v
                    +-------------+
                    |  constant   |
                    +-------------+
\end{verbatim}
\caption{関数呼び出し関係図}
\label{fig:callgraph}
\end{figure}

\subsection{関数の説明}
本節では，各関数がどのように文法規則に対応し，
入力をどのように消費するかを説明する。  
（ここに \verb|expression()|、\verb|term()|、\verb|factor()| の設計を追記）

%%%%%%%%%%%%%%%%%%%%%%%%%%%%%%%%%%%%%%%%%%%%
\section{工夫した点}
%%%%%%%%%%%%%%%%%%%%%%%%%%%%%%%%%%%%%%%%%%%%

（例）
\begin{itemize}
    \item 文法エラー位置を 1-based index で正確に表示できるように，
    \verb|next()| の設計を工夫した。
    \item 括弧・一元負号・指数など，
    優先順位の高い構文要素を自然に扱うため，関数分割を徹底した。
\end{itemize}

%%%%%%%%%%%%%%%%%%%%%%%%%%%%%%%%%%%%%%%%%%%%
\section{考察}
%%%%%%%%%%%%%%%%%%%%%%%%%%%%%%%%%%%%%%%%%%%%

（例）
\begin{itemize}
    \item 再帰下降構文解析は実装が簡潔で，
    構文規則の変更に強く拡張性が高い。
    \item 本課題では演算子優先順位が自然に表現される構造を実現した。
\end{itemize}

%%%%%%%%%%%%%%%%%%%%%%%%%%%%%%%%%%%%%%%%%%%%
\section{感想}
%%%%%%%%%%%%%%%%%%%%%%%%%%%%%%%%%%%%%%%%%%%%

（短めに）

%%%%%%%%%%%%%%%%%%%%%%%%%%%%%%%%%%%%%%%%%%%%
\section{作業工程}
%%%%%%%%%%%%%%%%%%%%%%%%%%%%%%%%%%%%%%%%%%%%

（例）
\begin{itemize}
    \item アルゴリズム設計：3 時間
    \item コーディング：6 時間
    \item デバッグ：3 時間
    \item レポート作成：3 時間
\end{itemize}

%%%%%%%%%%%%%%%%%%%%%%%%%%%%%%%%%%%%%%%%%%%%
\section{参考文献}
%%%%%%%%%%%%%%%%%%%%%%%%%%%%%%%%%%%%%%%%%%%%

\begin{itemize}
    \item プログラミングB「レポートの書き方」資料:contentReference[oaicite:1]{index=1}
    \item ChatGPT（使用部分を"*生成 AI*" として明示）
\end{itemize}

%%%%%%%%%%%%%%%%%%%%%%%%%%%%%%%%%%%%%%%%%%%%
\appendix
\section*{付録：プログラムリスト}

\begin{verbatim}
（ここに C プログラムを貼る）
\end{verbatim}

\end{document}
